\subsection{The block-size boundary and forest counterexamples}
\label{ian:gamma-set-blocks-and-forests}

Status: proved. Bounded block products satisfy
\(\Gamma^{\mathrm{set}}\) for block size at most two
(Lemma~\ref{ian:block-resampling-aa},
Corollary~\ref{ian:two-spin-blocks}), as do Ising forests with zero
edge bound (Proposition~\ref{ian:forest-addressing}), by the influence
bound of Corollary~\ref{ian:product-influence-aa-proved}, proved in
Subsection~\ref{ian:influence-form-uniform-aa}; these results no longer
assume (UAA). Block products fail
\(\Gamma^{\mathrm{set}}\) for every block-size bound at least three
(Proposition~\ref{ian:hard-block-addressing}), although
\(\Lambda_{\mathrm{blk}}(k)\) has hereditary full AA for every \(k\) by
Corollary~\ref{cor:resolution-product-alphabet-aa}; Ising forests fail
\(\Gamma^{\mathrm{set}}\) for every positive finite edge-strength bound
(Proposition~\ref{ian:forest-addressing}). The positive forest
simulation quoted there is
Theorem~\ref{thm:threshold-ising-forests}, whose bound is stated
through the uniform-input simulator \(s_U\) (via
\(S_B\) of Theorem~\ref{thm:threshold-block-products}); that simulation
form is supplied, with \(s_U(T,e)=O(T^6e^{-15})\), by
Corollary~\ref{cor:resolution-product-alphabet-aa} applied to the
uniform law.

The positive influence results below use Corollary~\ref{ian:product-influence-aa-proved} (Subsection~\ref{ian:influence-form-uniform-aa}), proved from (BC) and the OSSS covariance inequality with exponents \((12,12)\). An earlier version assumed instead the influence-form uniform theorem (UAA) imported in \path{influence/high_temperature/01_conventions_imported_uniform_theorem.tex}, Section~\ref{ht:conventions-and-imported-uniform-theorem}; that hypothesis is no longer used. Corollary~\ref{cor:resolution-product-alphabet-aa} supplies only simulation, and is used below only for simulation statements.

Let \(\Lambda_{\mathrm{blk}}(k)\) be the class of all products of arbitrary laws on disjoint coordinate blocks of size at most \(k\); zero atoms are allowed and block laws may depend on \(n\). This class is pinning-closed: conditioning
on a coordinate assignment conditions each block separately, leaving
independent residual blocks of size at most \(k\). Zero-mass
assignments are omitted.

For an independent block \(B\), define its resampling energy by
\[
 \mathcal E_B^\mu(f)
 :=\E_\mu\Var(f\mid X_{-B})
 =\frac12\E\bigl(f(X)-f(X_{-B},X'_B)\bigr)^2,
\]
where \(X'_B\) is an independent copy of the block.

\begin{lemma}[Block resampling AA without an atom floor]
\label{ian:block-resampling-aa}
Fix \(k\ge1\). There is an explicit
\(\kappa_k>0\) such that every multilinear polynomial \(f\) of degree at
most \(d\) with \(|f|\le1\) on the whole cube, on
every \(\mu\in\Lambda_{\mathrm{blk}}(k)\) satisfies
\[
 \Var_\mu f\ge v
 \quad\Longrightarrow\quad
 \max_B\mathcal E_B^\mu(f)\ge\kappa_k(v/d)^{12}.
\]
One choice is
\[
 \kappa_k=c_\star=\frac{9}{2^{61}\pi^{12}},
\]
the constant of Theorem~\ref{ian:uniform-influence-aa-proved}.
\end{lemma}

\begin{proof}
Regard each block \(B\) as one alphabet coordinate, with alphabet
\(\Omega_B\) the support of its block law \(\pi_B\) in
\(\{-1,1\}^B\). Zero-mass block states are simply deleted, so zero
atoms are allowed and no atom floor is needed. Then \(\mu\) is a
product law on \(\prod_B\Omega_B\). Every physical monomial of degree
at most \(d\) involves at most \(d\) blocks, so \(f\) has interaction
degree at most \(d\) on this product of block alphabets
(Definition~\ref{ian:alphabet-interaction-degree}), and it remains
bounded by one on the whole product set because it is bounded on the
entire physical cube. For the block coordinate \(B\), the quantity
\(\mathcal E_i^\mu\) of
Corollary~\ref{ian:product-influence-aa-proved} is exactly the
resampling energy \(\mathcal E_B^\mu(f)\). Assertion (A-AA) of that
corollary therefore gives
\[
 \max_B\mathcal E_B^\mu(f)
 \ge c_\star\frac{(\Var_\mu f)^{12}}{d^{12}}
 \ge\kappa_k(v/d)^{12}.
\]
\end{proof}

\begin{corollary}[Positive result through block size two]
\label{ian:two-spin-blocks}
Arbitrary coordinate-bias products satisfy
\[
 \max_iI_i^\mu(f)\ge c_\star(v/d)^{12},
\]
including deterministic coordinates. Products of arbitrary laws on
blocks of size at most two satisfy
\[
 \max_iI_i^\mu(f)\ge\frac{c_\star}{8}(v/d)^{12}.
\]
Thus both classes satisfy \(\Gamma^{\mathrm{set}}\) through the Local
alternative alone. The first bound extends
\eqref{ian:uniform-product-bound} from a common bias to
coordinate-dependent biases, at the same rate.
\end{corollary}

\begin{proof}
For a singleton block,
\[
 \mathcal E_{\{i\}}^\mu(f)
 =4p_i(1-p_i)I_i^\mu(f)\le I_i^\mu(f).
\]

For a block \(B\) of size at most two, fix all outside coordinates
and write the resulting function as \(h\). For any block states
\(u,w\),
\[
 (h(u)-h(w))^2
 \le4\sum_{i\in B}
       \bigl((D_ih(u))^2+(D_ih(w))^2\bigr).
 \tag{*}
\]
For adjacent states this is immediate. For opposite corners of a
square, telescope along each of the two length-two paths, apply
Cauchy--Schwarz on each path, and average the resulting inequalities.
This proves (*).

Integrating over independent block states and the outside coordinates
yields
\[
 \mathcal E_B^\mu(f)\le4\sum_{i\in B}I_i^\mu(f).
\]
Lemma~\ref{ian:block-resampling-aa} now proves both assertions.
No atom floor is used.
\end{proof}

\begin{remark}[Why size two escapes the size-three obstruction]
\label{ian:two-versus-three}
Two states of a block of size at most two are joined by a path of length at most two, and every edge of such a path touches one of its endpoints. Every derivative in (*) is therefore evaluated at a support state. In a block of size three, opposite corners are joined only through a middle edge whose endpoints may both have zero mass; the example of Proposition~\ref{ian:triple-random-set-failure} hides its variance exactly there.
\end{remark}

\begin{lemma}[Fixed-set reveal bound for addressing]
\label{ian:addressing-reveal}
Let independent components carry independent fair logical signs
\[
 A_1,\ldots,A_r,\qquad (Z_u)_{u\in\{-1,1\}^r}.
\]
Each logical sign is a measurable function of its own component.
Define
\[
 \operatorname{Addr}_r(A,Z)
 =\sum_{u\in\{-1,1\}^r}
       Z_u\prod_{j=1}^r\frac{1+u_jA_j}{2}.
\]
This function is Boolean and has variance one.

If a fixed physical coordinate set \(S\) touches \(s\) address
components and \(q\) data components, then
\[
 G(\operatorname{Addr}_r,S)
 \le q\,2^{s-2r}
 \le |S|\,2^{-r}.
\]
\end{lemma}

\begin{proof}
Reveal the entire components touched by \(S\); this can only increase
the reveal gain. Their additional information beyond their logical
signs does not affect the conditional expectation, by component
independence.

Let \(I\subseteq[r]\) be the touched address indices and
\(D\subseteq\{-1,1\}^r\) the touched data indices. The conditional mean
after this larger revelation is
\[
 2^{-(r-s)}
 \sum_{\substack{u\in D\\u|_I=A|_I}}Z_u.
\]
The revealed data signs are independent and centered. Each
\(u\in D\) is compatible with the revealed address signs with
probability \(2^{-s}\). The variance of this conditional mean is
therefore exactly
\[
 q\,2^{-s}2^{-2(r-s)}=q\,2^{s-2r}.
\]
Use \(s\le r\) and \(q\le|S|\).
\end{proof}

\begin{proposition}[Failure from block size three onward]
\label{ian:hard-block-addressing}
For every fixed \(k\ge3\),
\(\Gamma^{\mathrm{set}}(\Lambda_{\mathrm{blk}}(k))\) is false.
It already fails on products of uniform
\(\{(+,+,+),(-,-,-)\}\) blocks.
\end{proposition}

\begin{proof}
Take \(r+2^r\) independent three-spin blocks of this type. Label
\(r\) of them as address blocks and the others by
\(u\in\{-1,1\}^r\) as data blocks. In each block use its majority
\(M_3\) as the logical sign, and compose these signs with
\(\operatorname{Addr}_r\). Call the resulting physical polynomial
\(f_r\). It is multilinear and Boolean on the entire physical cube,
with
\[
 n_r=3(r+2^r),\qquad
 d_r:=3(r+1)\ge\deg f_r,\qquad
 \Var f_r=1.
\]

For a physical coordinate \(i\) in logical block \(B\), the
multilinear chain rule gives
\[
 D_if_r
 =(D_B\operatorname{Addr}_r)(\text{other logical signs})\,D_iM_3.
\]
The last factor vanishes on the block support. Thus
\[
 I_i(f_r)=0\quad\text{for every }i.
\]
Lemma~\ref{ian:addressing-reveal} gives, simultaneously for every
fixed physical set,
\[
 G(f_r,S)\le |S|\,2^{-r}.
\]

Fix any proposed Gamma constants and set \(v=1,d=d_r\). Let
\[
 L_r=\left\lceil B_Gd_r^{K_G}\right\rceil.
\]
For sufficiently large \(r\), \(n_r>L_r\) and
\[
 L_r2^{-r}<c_Gd_r^{-C_G}.
\]
Every allowed set fails the Global alternative, while every influence
is zero. This disproves Gamma for block size three, and hence for
every \(k\ge3\).
\end{proof}

\paragraph{Why block resampling does not finish the proposed route.}
Lemma~\ref{ian:block-resampling-aa} supplies a block with large
resampling energy, but that energy is not the variance revealed by
observing the block. Already \(f=Z_1Z_2\) has resampling energy one
in each block and zero singleton reveal gain.
The addressing construction makes the fixed-set obstruction
exponential in the degree.

The product-alphabet certificate and completion results in
\path{foundations/product_alphabets.tex},
Lemma~\ref{lem:resolution-alphabet-certificates} and
Corollary~\ref{cor:resolution-product-alphabet-aa}, are consistent
with this counterexample: they give adaptive simulation, which does
not require a small fixed revealing set. Indeed, the block-symbol
reduction in the proof of
Proposition~\ref{ian:triple-random-set-failure}, with three replaced
by \(k\), shows that \(\Lambda_{\mathrm{blk}}(k)\) has hereditary full
AA with \(O(kT^6e^{-15})\) physical queries for every \(k\); this
sharpens the bound of Theorem~\ref{thm:threshold-block-products} by
removing its block-length factor. With
Proposition~\ref{ian:hard-block-addressing} it recovers, for every
\(k\ge3\), that \(\Gamma^{\mathrm{set}}\) fails on a pinning-closed class
with hereditary full AA (for \(k=3\) this is
Proposition~\ref{ian:gamma-set-not-necessary}).

\begin{proposition}[Failure for Ising forests at every positive edge bound]
\label{ian:forest-addressing}
Fix \(0<\Theta=\beta J_{\mathrm{edge}}<\infty\).
The class of finite Ising forests with
\(|\beta J_{uv}|\le\Theta\), including arbitrary fields, fails
\(\Gamma^{\mathrm{set}}\). Zero fields and ferromagnetic couplings
already suffice.

For \(\Theta=0\), the class is a product class and satisfies Gamma,
by Corollary~\ref{ian:two-spin-blocks}.
\end{proposition}

\begin{proof}
Fix \(0<\theta\le\Theta\). For each odd \(r\ge3\), take
\(r+2^r\) disjoint stars, each with \(r\) leaves, zero fields, and
coupling \(\beta J=\theta\) on every edge.

For one star, its center \(C\) is fair, and conditional on \(C\)
the leaves are independent with
\[
 \Pr[X_i=C\mid C]=p:=\frac{1+\tanh\theta}{2}.
\]
Let \(M_r\) be the majority of its \(r\) leaves, represented by its
multilinear polynomial. It is Boolean on the entire cube, has degree
at most \(r\), and is fair by spin-flip symmetry.
Use these independent logical majorities as the address and data
signs in \(\operatorname{Addr}_r\), obtaining \(f_r\). Then
\[
 n_r=(r+1)(r+2^r),\qquad
 d_r:=r(r+1)\ge\deg f_r,\qquad
 |f_r|=1,\qquad \Var f_r=1.
\]

The polynomial does not depend on the centers, so their influences
are zero. For a leaf \(i\), write \(L_i\) for the set of leaves of its star;
\[
 D_iM_r=\mathbf1\left\{\sum_{j\in L_i\setminus\{i\}}X_j=0\right\}.
\]
Since the logical derivative of addressing has absolute value at
most one,
\[
 \begin{split}
 I_i(f_r)
 &\le\Pr\left[\sum_{j\in L_i\setminus\{i\}}X_j=0\right]\\
 &=\binom{r-1}{(r-1)/2}
        [p(1-p)]^{(r-1)/2}\\
 &\le(\operatorname{sech}\theta)^{r-1}.
 \end{split}
\]
The equality holds after conditioning on either center sign.

For any fixed physical set \(S\), reveal the full stars it touches.
Lemma~\ref{ian:addressing-reveal} then gives
\[
 G(f_r,S)\le |S|\,2^{-r}.
\]
For any proposed constants, put
\(L_r=\lceil B_Gd_r^{K_G}\rceil\).
Along sufficiently large odd \(r\),
\[
 (\operatorname{sech}\theta)^{r-1}<c_Ld_r^{-C_L},
 \qquad
 L_r2^{-r}<c_Gd_r^{-C_G},
 \qquad n_r>L_r.
\]
With \(v=1,d=d_r\), both alternatives fail.

Pinning preserves the forest class by deleting pinned vertices and
adding fields, without increasing the edge bound. The class
nevertheless has hereditary full AA by
\path{simulation/positive_examples/03_ising_forests.tex},
Theorem~\ref{thm:threshold-ising-forests}; that theorem's bound is
expressed through the uniform-input simulator \(s_U\), here supplied by
Corollary~\ref{cor:resolution-product-alphabet-aa} (see the Status
line). Thus fixed-set Gamma does not explain that positive simulation
theorem.

When \(\Theta=0\), the spins are independent, and
Corollary~\ref{ian:two-spin-blocks} applies.
\end{proof}
